\documentclass[conference]{IEEEtran}

% ── 中文与基础排版 ─────────────────────────────────────
\usepackage[UTF8]{ctex}
\usepackage{amsmath,amssymb}
\usepackage{booktabs}
\usepackage{graphicx}
\usepackage{url}
\usepackage{listings}
\usepackage{xcolor}

\lstset{
  basicstyle=\ttfamily\footnotesize,
  breaklines=true,
  frame=single,
  columns=fullflexible,
  keywordstyle=\color{blue},
  commentstyle=\color{gray}
}

\begin{document}

\title{面向课程知识库的检索增强生成系统设计与实现}

\author{
\IEEEauthorblockN{大数据学期项目方向 B}
\IEEEauthorblockA{RAG 检索增强生成助手\\
Python 3.11 + ChromaDB + Streamlit + AutoDL RTX 4090}
}

\maketitle

\begin{abstract}
本文设计并实现了一个面向课程知识库与技术问答资料的检索增强生成
（Retrieval-Augmented Generation, RAG）系统。系统覆盖文档摄取、
文本清洗、语义分块、结构化元数据提取、向量索引、混合检索、答案生成
与来源追溯等完整链路。针对本地 GTX 1660 SUPER 显存不足以稳定完成
百万级文本分块向量化的问题，本文采用 ``本地 ETL + 云端 GPU 推理''
的协同架构，在 AutoDL RTX 4090 上部署 OpenAI 兼容 Embedding 服务，
本地仅负责数据处理与 ChromaDB 写入。同时，系统新增 Bearer Token
保护远程推理接口，并对 Streamlit 前端输出进行 HTML 转义，降低公网
服务滥用与页面注入风险。最终项目形成 CLI 与 Web 双入口，配套 91 个
自动化测试用例，验证了摄取、预处理、检索、问答与安全渲染等核心行为。
\end{abstract}

\begin{IEEEkeywords}
RAG, ChromaDB, 向量检索, AutoDL, RTX 4090, Streamlit, 文档问答
\end{IEEEkeywords}

\section{引言}
课程项目、实验通知、知识点说明、Stack Overflow 问答和技术博客通常以
Markdown、PDF、HTML 片段等非结构化形式分散存在。传统关键词检索依赖
字面匹配，难以理解 ``交作业'' 与 ``提交方式''、``RAG'' 与
``检索增强生成'' 等语义等价表达。本文实现的 RAG 系统通过向量语义
召回和大语言模型答案生成，将杂乱资料转化为可追溯、可交互、可维护的
课程知识库助手。

本文的工程重点不只在问答效果，还包括可复现性与可运行性：第一，本地
显卡显存有限，百万级分块直接本地向量化会出现显存压力；第二，远程
GPU 服务若直接暴露到公网，存在被他人消耗算力的风险；第三，前端若
直接渲染模型输出，存在 HTML 注入隐患。因此，本文在系统架构中同时
加入云端算力卸载、安全鉴权和输出转义机制。

\section{系统总体架构}
系统采用分层流水线设计，如图式流程所示：

\begin{lstlisting}[language=bash,caption={RAG 数据流}]
data/raw + external corpus
  -> ingest.py        # md/txt/pdf/jsonl 摄取
  -> preprocess.py    # 清洗、分块、元数据提取
  -> embed_store.py   # local/remote/openai embedding
  -> ChromaDB         # HNSW + cosine
  -> query_parser.py  # 查询意图和过滤条件
  -> qa.py            # 答案生成与来源标注
  -> CLI / Streamlit
\end{lstlisting}

\begin{table}[!t]
\centering
\caption{核心模块职责}
\label{tab:modules}
\begin{tabular}{@{}ll@{}}
\toprule
\textbf{模块} & \textbf{职责} \\
\midrule
ingest.py & 读取 md/txt/pdf，JSONL 专用解析 \\
preprocess.py & 文本清洗、语义分块、元数据提取 \\
embed\_store.py & 嵌入生成、ChromaDB 写入与检索 \\
query\_parser.py & 自然语言查询解析为 search query 与 filters \\
qa.py & 基于检索片段生成答案并补全来源 \\
streamlit\_app.py & Web 交互、增量入库、调试状态展示 \\
embedding\_server.py & AutoDL 远程 OpenAI 兼容向量服务 \\
\bottomrule
\end{tabular}
\end{table}

\section{数据摄取与预处理}
\subsection{摄取边界}
文本摄取器只读取 \texttt{.md}、\texttt{.txt} 和 \texttt{.pdf}，
\texttt{.jsonl} 由专用函数读取。这样避免同一 JSONL 数据既被当作普通
文本摄取，又被结构化解析器再次摄取，导致重复入库。JSONL 记录会被映射
为统一结构：\texttt{source}、\texttt{path}、\texttt{text} 与
\texttt{fm\_meta}。

\subsection{清洗与分块}
系统通过正则流水线清除 HTML 标签、解码常见实体、删除控制字符并合并
空白行。分块采用段落优先、句子保护、贪心合并、长句滑窗的四层策略。
默认参数为 \texttt{chunk\_size=700}、\texttt{overlap=120}，在中文
课程资料中能兼顾语义完整性和向量聚焦度。

\subsection{元数据策略}
元数据策略支持 \texttt{merge}、\texttt{llm\_only} 与
\texttt{jsonl\_only}。最新版本对该参数进行白名单校验，非法值直接
抛出 \texttt{ValueError}，避免静默走错路径。Front-Matter 与 JSONL
元数据优先于 LLM 自动提取结果，以保证人工标注和原始数据字段的可信度。

\section{向量化与云端算力卸载}
\subsection{本地显存瓶颈}
项目初期尝试在本地 GTX 1660 SUPER 上完成大规模文本分块向量化。由于
本地显存较小，批量运行高维中文 Embedding 模型时容易触发显存不足或
吞吐不稳定。该问题不是单次查询阶段的瓶颈，而是集中建库阶段的典型
大数据工程瓶颈。

\subsection{AutoDL RTX 4090 远程服务}
为解决建库阶段算力问题，系统在 AutoDL RTX 4090 实例上部署
\texttt{scripts/embedding\_server.py}。该服务基于 FastAPI、PyTorch
和 SentenceTransformers，暴露 OpenAI 兼容的 \texttt{/v1/embeddings}
接口。本地 \texttt{embed\_store.py} 在 \texttt{OPENAI\_EMBEDDING\_MODEL=remote}
时，将文本批次发送给云端服务，并把返回的 1024 维向量写入本地 ChromaDB。

\begin{table}[!t]
\centering
\caption{Embedding 模式对比}
\label{tab:embedding}
\begin{tabular}{@{}llll@{}}
\toprule
\textbf{模式} & \textbf{模型} & \textbf{优势} & \textbf{适用场景} \\
\midrule
local & BAAI/bge-large-zh-v1.5 & 离线、零 API 成本 & 小批量或显存充足 \\
remote & BAAI/bge-large-zh-v1.5 & 4090 高吞吐 & 百万级建库 \\
openai-api & 兼容 Embedding API & 免维护 & 无本地/云端 GPU \\
\bottomrule
\end{tabular}
\end{table}

\subsection{公网服务安全}
AutoDL 自定义服务通常会获得公网地址。若不加限制，任何人都可能直接调用
Embedding 接口消耗 GPU 资源。因此系统新增 \texttt{EMBEDDING\_SERVER\_TOKEN}：
当云端配置该变量后，服务端会校验
\texttt{Authorization: Bearer <token>}；本地客户端读取同名变量并自动
作为 OpenAI-compatible API key 发送。未配置 token 时保留兼容模式，
便于课堂演示和本地调试。

\section{检索与答案生成}
\subsection{混合检索}
系统支持向量语义搜索、ChromaDB 元数据过滤和文档关键词过滤。对于多键
where 条件，系统将隐式多条件转换为显式 \texttt{\$and}，以适配 ChromaDB
查询语法。

\subsection{双路召回}
在未指定过滤条件时，系统采用双路召回：第一路优先检索课程相关分类
（wiki、faq、notice、case\_study、report、term、general），第二路
执行全库检索。两路结果按文本去重并按距离排序。若单路失败，系统记录
warning 并返回另一条成功路线；若两路均失败，则抛出 \texttt{RuntimeError}，
避免将配置错误误报为 ``未检索到内容''。

\subsection{答案生成与前端防注入}
答案生成模块要求 LLM 仅基于检索片段回答，并在回答中包含来源标注。
Streamlit 前端保留来源卡片的自定义 HTML，但模型回答、历史消息和错误
信息均先进行 HTML 转义，再转换换行为 \texttt{<br>}，防止知识库内容或
模型输出注入页面脚本。

\section{实现与接口配置}
\begin{table}[!t]
\centering
\caption{关键环境变量}
\label{tab:env}
\begin{tabular}{@{}ll@{}}
\toprule
\textbf{变量} & \textbf{用途} \\
\midrule
OPENAI\_API\_KEY & 对话模型 API Key \\
OPENAI\_MODEL & 答案生成与查询解析模型 \\
OPENAI\_EMBEDDING\_MODEL & local / remote / API 模型名 \\
OPENAI\_EMBEDDING\_BASE\_URL & 远程 Embedding 服务地址 \\
EMBEDDING\_SERVER\_TOKEN & 远程 Embedding 服务访问令牌 \\
LOCAL\_EMBEDDING\_MODEL & 本地或远程实际加载的模型名 \\
\bottomrule
\end{tabular}
\end{table}

典型远程配置如下：
\begin{lstlisting}[language=bash,caption={远程 4090 向量化配置}]
OPENAI_EMBEDDING_MODEL=remote
OPENAI_EMBEDDING_BASE_URL=https://<autodl-service>/v1
EMBEDDING_SERVER_TOKEN=<same-token-on-server>
LOCAL_EMBEDDING_MODEL=BAAI/bge-large-zh-v1.5
\end{lstlisting}

\section{测试与验证}
项目测试使用 \texttt{pytest}，通过 mock 隔离 OpenAI、ChromaDB 和外部网络。
当前共 91 个测试用例，覆盖摄取、预处理、向量检索、查询解析、问答生成、
Streamlit 安全渲染与集成链路。最新全量测试结果为：

\begin{lstlisting}[language=bash,caption={测试结果}]
python -m pytest tests/ -v
91 passed in 11.18s
\end{lstlisting}

\begin{table}[!t]
\centering
\caption{测试覆盖概览}
\label{tab:tests}
\begin{tabular}{@{}lc@{}}
\toprule
\textbf{模块} & \textbf{测试数} \\
\midrule
embed\_store & 15 \\
ingest & 27 \\
preprocess & 28 \\
qa & 9 \\
query\_parser & 5 \\
integration & 4 \\
app.rendering & 2 \\
\midrule
\textbf{总计} & \textbf{91} \\
\bottomrule
\end{tabular}
\end{table}

\section{成本与工程收益}
本项目将成本集中在建库阶段：远程 RTX 4090 负责高吞吐向量化，本地保留
ChromaDB 持久化结果用于后续演示和查询。这样既避开了本地显存不足，也
避免每次查询都依赖云端 GPU。与完全调用商业 Embedding API 相比，该
方案更适合一次性大规模建库、之后多次本地演示的课程项目场景。

\section{结论}
本文实现的 RAG 系统完成了从非结构化资料到可追溯问答的端到端闭环。
在工程上，系统重点解决了本地显存不足、远程服务公网暴露、前端模型输出
不可信、JSONL 重复摄取和检索错误静默降级等问题。最终系统以 AutoDL
RTX 4090 完成大规模向量化，以 ChromaDB 保存本地索引，以 Streamlit 和
CLI 提供双入口，并以 91 个自动化测试保证核心行为稳定。

\begin{thebibliography}{00}
\bibitem{rag}
P. Lewis et al., ``Retrieval-Augmented Generation for Knowledge-Intensive NLP Tasks,''
NeurIPS, 2020.
\bibitem{chroma}
Chroma, ``ChromaDB Documentation,'' \url{https://docs.trychroma.com/}.
\bibitem{bge}
BAAI, ``BGE Embedding Models,'' \url{https://huggingface.co/BAAI/bge-large-zh-v1.5}.
\bibitem{ieee}
IEEE, ``IEEE Conference Template,'' \url{https://www.ieee.org/conferences/publishing/templates.html}.
\end{thebibliography}

\end{document}
